\documentclass[11pt,a4paper]{article}
\usepackage[T1]{fontenc}
\usepackage{lmodern}
\usepackage[margin=27mm]{geometry}
\usepackage{amsmath,amssymb,amsthm,mathtools}
\usepackage{microtype,booktabs,longtable,array}
\usepackage{xcolor}
\usepackage[colorlinks=true,linkcolor=blue!45!black,citecolor=blue!45!black,urlcolor=blue!45!black]{hyperref}
\hypersetup{pdftitle={The exact quantum maximum of the five-outcome CGLMP inequality},pdfauthor={Dehao Lin},pdfsubject={Exact noncommutative SOS certificate; manuscript version 0.1.1}}
\pdfinfoomitdate=1
\pdftrailerid{}
\pdfsuppressptexinfo=15
\usepackage[numbers,sort&compress]{natbib}
\usepackage{fancyvrb}
\usepackage{enumitem}
\setlist{itemsep=2pt,topsep=4pt}
\setlength{\emergencystretch}{2em}
\numberwithin{equation}{section}
\newtheorem{theorem}{Theorem}
\newtheorem{lemma}[theorem]{Lemma}
\newtheorem{proposition}[theorem]{Proposition}
\newtheorem{corollary}[theorem]{Corollary}
\theoremstyle{remark}\newtheorem{remark}[theorem]{Remark}
\newcommand{\C}{\mathbb C}
\newcommand{\Q}{\mathbb Q}
\newcommand{\HH}{\mathcal H}
\newcommand{\id}{\mathbf 1}
\newcommand{\tr}{\operatorname{tr}}
\newcommand{\NF}{\operatorname{NF}}
\newcommand{\diag}{\operatorname{diag}}
\newcommand{\ket}[1]{\lvert #1\rangle}
\newcommand{\bra}[1]{\langle #1\rvert}
\newcommand{\adj}{\dagger}
\title{The exact quantum maximum of the\newline five-outcome CGLMP inequality}
\author{Dehao Lin\\[3pt]
  \small School of Physics, Sun Yat-sen University, Guangzhou, China\\
  \small ORCID: \href{https://orcid.org/0009-0001-4551-8490}{0009-0001-4551-8490}}
\date{4 October 2026\quad---\quad Manuscript version 0.1.1}
\begin{document}
\maketitle
\begin{abstract}
We give a computer-assisted exact certificate for the quantum maximum of the
\mbox{standard} five-outcome Collins--Gisin--Linden--Massar--Popescu (CGLMP) inequality,
in the normalization with local bound $2$. The maximum is the largest real root
$\mu$ of $5\mu^6-65\mu^4+144\mu^2+96\mu+16=0$, namely
$3.0157104755226732855\ldots$. An explicit noncommutative sum of squares proves
the upper bound for arbitrary local Hilbert spaces and arbitrary five-outcome
local positive-operator-valued measurements. The proof first establishes a
bounded commuting-projector operator inequality and then uses a fixed common
embedding for both settings of each local measurement. A nonmaximally entangled
state in $\C^5\otimes\C^5$, with standard Fourier measurements, attains the
bound. We give all exact coefficient data through five small Hermitian Gram
blocks, distinguish quotient-ring identities from positivity in the selected
embedding, and specify finite rational-arithmetic verification procedures.
The result does not require convergence of a numerical optimizer or of a
semidefinite hierarchy. It is a computer-assisted proof, not a proof-assistant
formalization; no claim of uniqueness, self-testing, or absolute priority is made.
\end{abstract}

\section{Introduction and relation to prior work}
Bell inequalities compare correlations obtainable from local classical models
with those obtainable from quantum measurements. The CGLMP family provides a
two-party, two-setting inequality for each finite number of outcomes
\cite{CGLMP}. The number of outcomes and the dimension of a local quantum system
are different parameters: establishing an optimum within a prescribed
five-dimensional family does not, by itself, give a universal five-outcome
quantum bound.

The standard Fourier measurement family and the improvement obtained with
nonmaximally entangled states have a substantial history. Zohren and Gill
\cite{ZohrenGill} studied the optimal-state structure and large-outcome
behavior numerically. Hu, Deng, and Chen \cite{HuDengChen} investigated positive
and negative violations by a Bell-operator method. Their numerical values and
measurement candidates are background for the present work, rather than
premises for its upper bound. In particular, the candidate decimal value
reported here is not a newly discovered numerical prediction.

The Navascu\'es--Pironio--Ac\'in hierarchy provides semidefinite upper bounds
on quantum correlations \cite{NPA2007,NPA2008}. Its dual perspective suggests
noncommutative sum-of-squares (SOS) certificates. Once an exact positive
certificate has been specified, its soundness can be checked directly, without
rerunning the optimization that helped find it. Ioannou and Rosset
\cite{IoannouRosset} developed symmetry-adapted methods and gave exact CGLMP
certificates for three and four outcomes. Their version 2, Section VIII.4,
reports a numerical five-outcome calculation and explicitly stops before an
exact SOS decomposition; Section VII already describes exact candidate
bounds and states. This comparison is specific to that version and those
claims, and is not an assertion of exhaustive priority clearance.

The object supplied here is an exact five-outcome positive certificate with
its actual algebraic embedding and a complete analytic connection to the
local-POVM model. Its mathematical content is unchanged from the frozen
scientific payload in \cite{Artifact}. This manuscript makes the coefficient
data, root selection, positivity, physical attainment, and infinite-dimensional
interfaces explicit in one place. The appendices contain the finite data
needed to reconstruct the certificate without a numerical search.

\section{Scenario, conventions, and theorem}\label{sec:theorem}
Let $\HH_A,\HH_B$ be arbitrary complex Hilbert spaces, with no finite-dimensional
or separability assumption. For $x,y\in\{0,1\}$ let
$\{E_{a|x}\}_{a=0}^4$ and $\{F_{b|y}\}_{b=0}^4$ be local POVMs:
$E_{a|x},F_{b|y}\succeq0$ and each set sums to the identity on its own space.
For a density operator $\rho$ on $\HH_A\otimes\HH_B$, define
\begin{equation}
 P(a,b|x,y)=\tr\!\left[\rho(E_{a|x}\otimes F_{b|y})\right].
\end{equation}
The argument also holds for any normalized positive functional on the
bounded-operator algebra, since the upper bound is an operator inequality.
The symbols $A_x,B_y$ below denote measurement outcomes, not Hermitian
observables. All outcome equalities are modulo $5$.

The functional is
\begin{equation}\label{eq:cglmp}
\begin{aligned}
 I_5=\sum_{k=0}^{1}\left(1-\frac{k}{2}\right)\big[&
 P(A_0=B_0+k)+P(B_0=A_1+k+1)\\
 &+P(A_1=B_1+k)+P(B_1=A_0+k)\\
 &-P(A_0=B_0-k-1)-P(B_0=A_1-k)\\
 &-P(A_1=B_1-k-1)-P(B_1=A_0-k-1)\big].
\end{aligned}
\end{equation}
For example, $P(A_x=B_y+k)=\sum_{b=0}^4P(b+k,b|x,y)$.

\begin{theorem}\label{thm:main}
Let $\mu$ be the largest real root of
\begin{equation}\label{eq:sextic}
 p(t)=5t^6-65t^4+144t^2+96t+16.
\end{equation}
Every admissible strategy on arbitrary local Hilbert spaces, with an
arbitrary normalized bipartite state and arbitrary five-outcome local POVMs,
satisfies $I_5\leq\mu$. Taking the supremum over all such local spaces,
states, and POVMs gives
\begin{equation}\label{eq:maximum}
 \sup I_5=\mu=3.015710475522673285523983801961\ldots.
\end{equation}
The supremum is attained by the explicit projective strategy in
Section~\ref{sec:attainment} on $\C^5\otimes\C^5$. The local hidden-variable
bound in this normalization is $2$.
\end{theorem}

For later use, put $[z]_5\in\{0,1,2,3,4\}$, $f(z)=1-z/2$, and relabel
\begin{equation}\label{eq:relabel}
 D_0=A_0,\quad D_1=B_0,\quad D_2=A_1+1,\quad D_3=B_1+1.
\end{equation}
Directly comparing each pairwise event coefficient gives
\begin{equation}\label{eq:cyclic}
 I_5=\mathbb E\left[
 \sum_{r=0}^{2}f([D_r-D_{r+1}]_5)+f([D_3-D_0-1]_5)\right].
\end{equation}
This notation is a sum of expectations from four experiments; it does not
postulate a joint distribution for incompatible settings. For a deterministic
local assignment, the four residues sum to a nonnegative integer congruent
to $4$ modulo $5$, and hence to at least $4$. Thus $I_5\leq4-4/2=2$.
The assignment $A_0=A_1=B_0=B_1=0$ attains $2$, and convexity proves the
local bound for arbitrary local randomization.

\section{Bell polynomial and algebraic embedding}\label{sec:algebra}
\subsection{The universal projective algebra}
Temporarily assume the measurements are projective. Let
$\omega=e^{2\pi i/5}$ and encode the relabelled outcomes by unitaries
\begin{equation}
 U_r=\sum_{a=0}^4\omega^aP^{(r)}_a,\qquad r=0,1,2,3.
\end{equation}
The only relations used in the upper bound are
\begin{equation}\label{eq:relations}
 U_r^5=\id,\quad U_r^\adj=U_r^{-1},\quad
 [U_r,U_s]=0\quad(r\ \text{even},\ s\ \text{odd}).
\end{equation}
There is no same-party commutation assumption.
The discrete Fourier coefficients of $f$ are
\begin{equation}\label{eq:fourier}
 c_0=0,\qquad
 c_k=\frac15\sum_{z=0}^4\left(1-\frac z2\right)\omega^{-kz}
     =\frac1{2(1-\omega^{-k})},\quad 1\leq k\leq4.
\end{equation}
Fourier synthesis in \eqref{eq:cyclic} yields the self-adjoint Bell polynomial
\begin{equation}\label{eq:bell}
 B=\sum_{r=0}^2\sum_{k=1}^4c_kU_r^kU_{r+1}^{-k}
   +\sum_{k=1}^4c_k\omega^{-k}U_3^kU_0^{-k}.
\end{equation}
In every projective realization its expectation is exactly \eqref{eq:cglmp}.
For exact checking, the finite sum in \eqref{eq:fourier} avoids algebraic
division and independently fixes the normalization and wraparound sign.

\subsection{The selected actual root}
We have $p(3)=-20$ and $p(31/10)>0$, while
\begin{equation}
 p'(t)=t^3(30t^2-260)+288t+96>0\qquad(t\geq3).
\end{equation}
Thus $(3,31/10)$ contains exactly one root, and no larger real root exists.
Set
\begin{equation}\label{eq:tower}
 s=\sqrt5>0,\qquad u=\sqrt{10+2s}>0,\qquad x=5\mu.
\end{equation}
The polynomial identity
\begin{equation}
 [5(t^3-4t-4)]^2-5(5t^2-4t-8)^2=5p(t)
\end{equation}
selects the positive square-root branch, since both bracketed quantities are
positive for $t\geq3$. Consequently
\begin{align}
 \mu^3&=s\mu^2+(4-4s/5)\mu+4-8s/5,\label{eq:cubicmu}\\
 x^3&=5sx^2+(100-20s)x+500-200s.\label{eq:cubic}
\end{align}
Define
\begin{equation}\label{eq:zeta}
 \zeta=\frac{u+i(s-1)}4.
\end{equation}
The relations in \eqref{eq:tower} give $|\zeta|=1$, $\zeta^5=i$,
$\Im\zeta>0$, and $\Re\zeta>1/\sqrt2$. Among the fifth roots of $i$ these
conditions select $\zeta=e^{i\pi/10}$, so that $\omega=\zeta^4$.

\subsection{Formal identities versus actual positivity}\label{sec:ring}
The serialized scalar convention is
\begin{equation}\label{eq:scalar}
 \langle n;d\rangle=
 \frac1d\sum_{\substack{a,c,e\in\{0,1\}\\b\in\{0,1,2\}}}
 n_{a+2b+6c+12e}\,s^ax^bu^ci^e,\qquad d\in\mathbb Z_{>0}.
\end{equation}
Every listed numerator is an integer. Products can be reduced with
$s^2=5$, $u^2=10+2s$, $i^2=-1$, and \eqref{eq:cubic}.
This is a quotient-ring calculation. It need not assume that these relations
give an irreducible field presentation: a zero obtained using them evaluates
to zero at the specified actual complex numbers. Conversely, a formal zero
does not select a real root or prove that a scalar is positive.

For a separate implementation one may use the 24 monomials $\zeta^ax^b$,
$0\leq a<8$, $0\leq b<3$, and
\begin{equation}\label{eq:cyclotomic}
 \Phi_{20}(z)=z^8-z^6+z^4-z^2+1,
 \quad s=2(\zeta^2+\zeta^{-2})-1,\quad u=2(\zeta+\zeta^{-1}),\quad i=\zeta^5.
\end{equation}
The same soundness statement applies. In particular, compatibility of a
different embedding with the formal identity does not imply positivity in
that embedding. Positivity below is proved only after the actual embedding
\eqref{eq:tower}--\eqref{eq:zeta} has been fixed.

\section{An explicit finite sum-of-squares certificate}\label{sec:sos}
\subsection{Generating words and exact Gram data}
Index from zero the ordered lists
\begin{align}
 a&=(\id,U_0,U_0^2,U_0^3,U_0^4,U_2,U_2^2,U_2^3,U_2^4),\\
 b&=(\id,U_1,U_1^2,U_1^3,U_1^4,U_3,U_3^2,U_3^3,U_3^4),
 \qquad q_{9r+t}=a_rb_t.\label{eq:q}
\end{align}
Thus $q$ has 81 entries with an unambiguous tensor-list order.
There are five blocks, labelled $k=1,6,7,8,9$, of sizes
$n_k=2,4,3,3,2$. Their labels are identifiers, not consecutive indices.

Appendix~\ref{app:data} specifies $81\times4$ phase matrices $F_k$,
$4\times n_k$ matrices $K_k$, and Hermitian $n_k\times n_k$ matrices $H_k$.
Set $E_k=F_kK_k$ and
\begin{equation}\label{eq:gram}
 G=\sum_{k\in\{1,6,7,8,9\}}E_kH_kE_k^\adj.
\end{equation}
The exact coefficient identity is
\begin{equation}\label{eq:gramidentity}
 \mu\id-B=\sum_{i,j=0}^{80}
 \left(G_{ij}q_i^\adj q_j+\overline{G_{ij}}q_iq_j^\adj\right).
\end{equation}
All of $F_k,K_k,H_k$ are printed in the appendix, including the integer
numerators and denominators. Equation~\eqref{eq:gramidentity} therefore
specifies a finite, reproducible algebraic claim, rather than an instruction
to find a feasible matrix.

\begin{lemma}[Exact certificate]\label{lem:certificate}
For the data of Appendix~\ref{app:data}, every $H_k$ is positive definite in
the actual embedding \eqref{eq:tower}, and \eqref{eq:gramidentity} holds
modulo \eqref{eq:relations}. Equivalently, there are exactly 14 positive
weights and 14 explicitly determined polynomials such that
\begin{equation}\label{eq:sos}
 \mu\id-B=\sum_{k}\sum_{r=0}^{n_k-1}d_{kr}
 (R_{kr}^\adj R_{kr}+R_{kr}R_{kr}^\adj),\qquad d_{kr}>0.
\end{equation}
These are the same weights and polynomials as in the frozen compact
certificate \texttt{SOS14.json}.
\end{lemma}

\subsection{Exact identity check}\label{sec:identity}
Here is a terminating coefficient-wise verification of the algebraic part.
Represent a word by an ordered list $(r_1,m_1),\ldots,(r_\ell,m_\ell)$,
meaning $U_{r_1}^{m_1}\cdots U_{r_\ell}^{m_\ell}$. Move even-index factors
to the left and odd-index factors to the right, preserving their relative
order within each party. In each of the two lists, combine adjacent powers
of the same generator modulo $5$, deleting powers congruent to zero and
repeating if a new adjacency is created. The result is a reduced word in
$(C_5*C_5)\times(C_5*C_5)$. This normal form is unique. The involution
reverses the list and negates powers before reduction.

Form $G$ by exact scalar arithmetic and expand the right-hand side of
\eqref{eq:gramidentity}. For each normal word $w$, collect the coefficient
\begin{equation}\label{eq:wordtest}
 \sum_{i,j:\,\NF(q_i^\adj q_j)=w}G_{ij}
 +\sum_{i,j:\,\NF(q_iq_j^\adj)=w}\overline{G_{ij}}
 -[w](\mu\id-B).
\end{equation}
Reducing its scalar coefficients by \eqref{eq:cubic} and
\eqref{eq:cyclotomic} gives zero for all 1,681 distinct words generated
by the full 81-word sequence. Words outside that universe occur on neither
side. This computation never commutes $U_0$ with $U_2$, or $U_1$ with $U_3$.
The supplied verifier uses the finite Fourier sum in \eqref{eq:fourier}
to construct its target independently of any stored target string.

For clarity, the algebraic statement can also be tested directly after
forming \eqref{eq:sos}: its 14 polynomials have 164 nonzero coefficients in
total, and the union of expanded support with the Bell target has 273
normal words. The counts 273 and 1,681 refer to different representations
of the same identity, not to incompatible universes of verification.

\subsection{Strict positivity with rational enclosures}\label{sec:positive}
Let $\Delta_{kj}$ be the determinant of the leading $j\times j$ submatrix
of $H_k$, with $\Delta_{k0}=1$. Determinants involve only additions and
multiplications, so their exact computation does not presume that any pivot
is nonzero. Appendix~\ref{app:intervals} gives rational boxes for $\mu,s,u$
and strict rational enclosures for every $\Delta_{kj}$. All lower endpoints
are positive. Each determinant is checked to be exactly real by conjugation.
Sylvester's criterion proves that each $H_k$ is positive definite.

For completeness, interval arithmetic here has its elementary inclusion
meaning: $[a,b]+[c,d]=[a+c,b+d]$, and interval multiplication takes the minimum
and maximum of $ac,ad,bc,bd$. Scalar signs are handled by taking the smaller
and larger scaled endpoints. Substitute the checked boxes into either the
real tower monomials or the real and imaginary parts of the cyclotomic
monomials, and sum the exact rational multiples. No decimal rounding enters
the test. The boxes themselves are checked by $p(\mu_-)<0<p(\mu_+)$,
$s_-^2<5<s_+^2$, and
$u_-^2<10+2s_-\leq10+2s_+<u_+^2$, with positive endpoints.

The positive leading minors now justify the unit-lower-triangular
factorization $H_k=L_kD_kL_k^\adj$. It is determined without choices by
\begin{align}
 d_{kr}&=(H_k)_{rr}-\sum_{t<r}|(L_k)_{rt}|^2d_{kt},\label{eq:ldl1}\\
 (L_k)_{jr}&=\frac{(H_k)_{jr}-\sum_{t<r}(L_k)_{jt}d_{kt}
                       \overline{(L_k)_{rt}}}{d_{kr}}\quad(j>r),\label{eq:ldl2}
\end{align}
with $(L_k)_{rr}=1$ and $(L_k)_{jr}=0$ for $j<r$.
Indeed $d_{kr}=\Delta_{k,r+1}/\Delta_{kr}>0$.
Define
\begin{equation}\label{eq:R}
 R_{kr}=\sum_{i=0}^{80}\overline{(E_kL_k)_{ir}}\,q_i.
\end{equation}
Expanding \eqref{eq:R} gives \eqref{eq:sos} from \eqref{eq:gramidentity}.
The conjugation in this formula is important: it fixes the orientation
of the Gram convention. Exact comparison with the frozen compact table
checks all $81\cdot14=1{,}134$ coefficient positions, including zeros, and
all 14 weights. The smallest weight is the one labelled $9{:}1$ and lies
strictly between $3584887/10^{10}$ and $3584888/10^{10}$.

\begin{remark}
The five small blocks are positive definite. The $81\times81$ Gram matrix
$G$, whose rank is at most 14, is only claimed to be positive semidefinite.
The printed Gram data and the serialized SOS data are two linked
presentations; agreement of one presentation with its target does not, by
itself, authenticate every byte of the other.
\end{remark}

\section{From the certificate to arbitrary local POVMs}\label{sec:bridge}
\subsection{Bounded representations}
Evaluate \eqref{eq:sos} in any bounded-operator representation of
\eqref{eq:relations} on any Hilbert space. For every vector $\xi$,
\begin{equation}\label{eq:positiverepresentation}
 \langle\xi,(\mu\id-B)\xi\rangle
 =\sum_{k,r}d_{kr}\left(\|R_{kr}\xi\|^2+
                                    \|R_{kr}^\adj\xi\|^2\right)\geq0.
\end{equation}
The sum is finite and its factors are bounded polynomials in unitaries.
There is no truncation, domain problem, or interchange of limits.
Thus $B\preceq\mu\id$ for arbitrary commuting projective representations,
and in particular for all tensor-product projective strategies.

\subsection{A common embedding for both settings}\label{sec:dilation}
We give an explicit local dilation to avoid any setting-dependent state
embedding. On one party's arbitrary Hilbert space $\HH$, let
$\{M_{a|x}\}_{a=0}^4$ be the POVM for setting $x=0,1$. Put
$\mathcal L=\HH^{\oplus5}$, with coordinate projections $\Pi_a$, and define
\begin{equation}
 V_x:\HH\longrightarrow\mathcal L,\qquad
 V_xh=(\sqrt{M_{0|x}}h,\ldots,\sqrt{M_{4|x}}h).
\end{equation}
Positivity gives bounded positive square roots and
$V_x^\adj V_x=\sum_aM_{a|x}=\id$. The defect
$D_x=\id_{\mathcal L}-V_xV_x^\adj$ is an orthogonal projection. On the
\emph{same} space $\widetilde\HH=\HH\oplus\mathcal L$ for both settings set
\begin{equation}\label{eq:fixedJ}
 J h=(h,0),\qquad
 W_x=\begin{pmatrix}0&V_x^\adj\\ V_x&D_x\end{pmatrix},\qquad
 Q_a=\begin{pmatrix}\delta_{a0}\id_{\HH}&0\\0&\Pi_a\end{pmatrix}.
\end{equation}
The $Q_a$ are five mutually orthogonal projections summing to the identity;
their upper block assigns the unused summand to outcome $0$. Moreover,
\begin{equation}
 W_x^\adj=W_x,\qquad
 W_x^2=\begin{pmatrix}
 V_x^\adj V_x&V_x^\adj D_x\\D_xV_x&V_xV_x^\adj+D_x^2
 \end{pmatrix}=\id.
\end{equation}
This uses $V_x^\adj V_x=\id$, not the generally false assertion
$V_xV_x^\adj=\id$. It holds also for proper isometries in infinite dimension.
The projectors $\widetilde M_{a|x}=W_xQ_aW_x$ therefore satisfy
\begin{equation}\label{eq:compression}
 J^\adj\widetilde M_{a|x}J=V_x^\adj\Pi_aV_x=M_{a|x}
 \quad\text{for every }a,x,
\end{equation}
with precisely the same fixed $J$ in both equations for $x$.

Apply this construction separately to Alice and Bob. With
$\mathcal J=J_A\otimes J_B$ we have, for every $a,b,x,y$,
\begin{equation}\label{eq:jointcompression}
 \mathcal J^\adj(\widetilde E_{a|x}\otimes\widetilde F_{b|y})\mathcal J
 =E_{a|x}\otimes F_{b|y}.
\end{equation}
Thus the single embedded state $\widetilde\rho=\mathcal J\rho\mathcal J^\adj$
preserves all joint probabilities simultaneously. Zero effects and mixed
states cause no difficulty. For a general positive normalized functional
$\varphi$, the formula $T\mapsto\varphi(\mathcal J^\adj T\mathcal J)$ gives
the corresponding embedded state functional.

The local tensor-product construction preserves cross-party commutation, so
the projective upper bound gives $I_5\leq\mu$ for every POVM strategy in
Theorem~\ref{thm:main}. Compression has only been used in
\eqref{eq:jointcompression}, where the two local factors separate. We have
not treated compression as a homomorphism on arbitrary products of
same-party operators. Nor have we asserted a simultaneous dilation theorem
for arbitrary abstract commuting POVMs: the POVM scope here is the stated
local tensor-product model.

\section{An explicit attaining strategy}\label{sec:attainment}
Let $\ket j$, $0\leq j\leq4$, be the standard basis of $\C^5$.
Use the projectors onto
\begin{align}
 \ket{a;x}&=\frac1{\sqrt5}\sum_{j=0}^4
 e^{2\pi i j(a+\alpha_x)/5}\ket j,
 &&(\alpha_0,\alpha_1)=(0,1/2),\label{eq:alice}\\
 \ket{b;y}&=\frac1{\sqrt5}\sum_{j=0}^4
 e^{2\pi i j(-b+\beta_y)/5}\ket j,
 &&(\beta_0,\beta_1)=(1/4,-1/4).\label{eq:bob}
\end{align}
For each setting these vectors form an orthonormal basis by the elementary
sum of fifth roots of unity. Define
\begin{equation}\label{eq:fs}
 t=\frac u2,\quad f_1=\frac1t,\quad f_2=\frac{s-1}2,\quad
 f_3=\frac ts,\quad f_4=\frac{s+1}2,
\end{equation}
and
\begin{equation}\label{eq:ab}
 a=\frac{\mu(f_1+f_3)+2f_1f_2}
          {\mu(\mu-f_2)-2f_1^2},\qquad
 b=\frac{2(f_2+f_1a)}\mu.
\end{equation}
Since $\mu>3$ and $0<f_1,f_2<1$, the denominator defining $a$ exceeds $4$.
All divisions are legitimate and $a,b>0$. Numerically,
$a=0.7188787436302164244\ldots$ and
$b=0.6605215707550941694\ldots$.
The state is
\begin{equation}\label{eq:state}
 \ket\psi=\frac{\ket{00}+a\ket{11}+b\ket{22}+a\ket{33}+\ket{44}}
                    {\sqrt{2+2a^2+b^2}}.
\end{equation}

We spell out attainment analytically. In an equality event, summing over
the common outcome in the Fourier projectors kills matrix elements from
$\ket{ll}$ to $\ket{jm}$ unless $j=m\pmod5$. Consequently the subspace
$\operatorname{span}\{\ket{jj}:0\leq j\leq4\}$ is invariant under $B$.
Substitution of \eqref{eq:alice}--\eqref{eq:bob} into
\eqref{eq:cglmp} gives on this subspace
\begin{equation}\label{eq:Toeplitz}
 T_{jj}=0,\qquad
 T_{jl}=\frac1{2\cos(\pi(j-l)/10)}=f_{|j-l|}\quad(j\ne l).
\end{equation}
This follows by summing the four setting-pair phases; equivalently it is
the five-point Fourier sum \eqref{eq:fourier} evaluated with the stated
shifts. The trigonometric values in \eqref{eq:Toeplitz} are exactly those
in \eqref{eq:fs}.

For $\gamma=(1,a,b,a,1)^{\mathsf T}$, the independent row equations are
\begin{align}
 \mu&=f_4+(f_1+f_3)a+f_2b,\label{eq:row0}\\
 \mu a&=f_1+f_3+f_2a+f_1b,\label{eq:row1}\\
 \mu b&=2f_2+2f_1a.\label{eq:row2}
\end{align}
Equations~\eqref{eq:row1}--\eqref{eq:row2} give \eqref{eq:ab}.
The determinant of the reduced three-dimensional eigenvalue equation is
\begin{equation}
 \det\begin{pmatrix}
 z-f_4&-(f_1+f_3)&-f_2\\
 -(f_1+f_3)&z-f_2&-f_1\\
 -2f_2&-2f_1&z
 \end{pmatrix}
 =z^3-sz^2-(4-4s/5)z-4+8s/5.
\end{equation}
At $z=\mu$ this vanishes by \eqref{eq:cubicmu}; the nonzero denominator
already established implies \eqref{eq:row0}. Hence $T\gamma=\mu\gamma$.
Together with invariance this proves $B\ket\psi=\mu\ket\psi$ in the full
25-dimensional space. It proves $I_5=\mu$ and completes the proof of
Theorem~\ref{thm:main}.

The artifact additionally reconstructs all 100 projector-to-unitary matrix
entries and checks the 25 full-space eigenvector coordinates exactly. These
checks guard against phase, outcome-label, and tensor-order transcription
errors; they do not replace the analytic universal upper-bound argument.

\section{Computer-assisted verification and reproducibility}\label{sec:repro}
\subsection{What is finite, and what is analytic}
The load-bearing finite claims are the exact coefficient identity,
Hermiticity and strict positivity in the selected embedding, and the exact
attaining strategy. The bounded-representation implication and the
common-embedding local-POVM bridge are analytic arguments given above.
There is no inference from a finite sample of dimensions to arbitrary
Hilbert spaces.

The discovery calculation used symmetry reduction, numerical exploration,
and exact reconstruction. Five retained blocks produced 42 real parameters;
the exact linear system had rank 34 and eight free parameters were fixed
rationally. This history motivates the data but is not an assumption of
the final proof. Appendix~\ref{app:data} supplies the resulting parameters
directly; no search, solver status, numerical matrix-rank guess, or
hierarchy-convergence argument is needed for verification.

\subsection{Replay paths and their dependence}
The frozen payload provides three relevant entry points:
\begin{itemize}
\item \texttt{verify\_independent.py}: full 1,681-word Gram identity,
exact LDL reconstruction, rational positivity, and physical attainment;
\item \texttt{verify\_sos14.py}: direct expansion of the 14 compact
anticommutators, without reading Gram data;
\item \texttt{verify\_statement.py}: the root bridge, stated Schmidt
amplitudes, denominator bound, reduced eigenvector, and Fourier
orthonormality. The literal event convention is checked by the physical
routine used in the first two entry points.
\end{itemize}
All three paths share scalar arithmetic in the distributed implementation;
they are not separate independent votes. The name of an entry point does not
establish independence of its author, runtime, or mathematical premises.
The release-level hardened entry points add strict schema and target/input
binding, without changing the frozen numerical coefficients.

The paper's \texttt{check\_paper\_data.py} regenerates the printed data,
checks the phase-factorization $E_k=F_kK_k$, reconstructs every Hermitian
matrix, and verifies agreement with all canonical compact coefficients and
weights. It computes the leading-minor intervals in
Appendix~\ref{app:intervals} by determinant expansion and rational interval
arithmetic. This is a data-correspondence check using the published exact
arithmetic implementation, not a claim of a new independent implementation.

The repository preserves the frozen payload, historical receipts, and later
adversarial-audit evidence separately. Subsequent multi-audit work tested
literal events, alternate coefficient representations, same-party word
order, actual branch signs, generic operator representations, and the fixed
POVM embedding. Failed and partial audit runs are retained and distinguished
from successful replays. Finite-field and floating tests are supplementary:
they prove neither characteristic-zero identity nor real positivity.
Agreement among audits is not itself an argument for the theorem.

\subsection{Practical reproducibility}
The four principal frozen certificate files are
\path{EXACT_KERNELS.json}, \path{EXACT_SOS_CANDIDATE.json},
\path{POSITIVITY_CERTIFICATE.json}, and \path{SOS14.json}.
Their release manifests identify exact bytes. Descriptive discovery status
strings, cached pass flags, decimal spectra, and historical file pointers are
not mathematical premises. A successful replay must check its intended
semantic result, not merely a zero exit code. The release records sources,
inputs, tool versions, and outputs, and rejects unsupported optimized-Python
execution where assertions could otherwise be stripped.

The manuscript is buildable with standard pdfLaTeX and BibTeX; the source,
complete coefficient appendix, bibliography, and clean build command are
distributed with it. The full scientific artifact is available at
\url{https://github.com/LeoLam233/cglmp5-exact-value} \cite{Artifact}.
Hashes are provenance aids: they identify data and code but do not prove
the truth of a mathematical assertion or certify an unobserved execution.

\section{Limitations and authorship disclosure}\label{sec:limits}
This theorem concerns exactly five outcomes per setting. It does not
establish a formula for all outcome numbers, uniqueness of the realizing
measurements or state, a self-testing theorem, or an optimal statistical
strength claim. The attaining state is nonmaximally entangled, but its
existence does not classify every maximizer. The commuting-projector upper
bound and the local tensor-product POVM theorem have been deliberately
distinguished.

The proof depends on the correctness of explicit finite exact calculations
and their elementary analytic interfaces. Public data and multiple
reconstructions reduce specific risks; no finite testing campaign eliminates
all possible software, runtime, or hardware errors. Neither this release nor
the accompanying audit verdict is human peer review, formal proof-assistant
verification, or an absolute originality/priority guarantee.

\paragraph{AI assistance and responsibility.}
Dehao Lin is the sole human author. Frontier AI systems substantially assisted
with problem solving, proof construction, exact-certificate generation,
verification software, adversarial auditing, literature triage, manuscript
preparation, and release work. No AI system is a manuscript author. The human
author selected and supervised the research workflow, decided what to
preserve and publish, and takes responsibility for the public claims and
artifact curation. Repository contributor attribution is distinct from
academic authorship.

\bibliographystyle{unsrtnat}
\bibliography{references}

\clearpage
\appendix
\section{Complete exact coefficient data}\label{app:data}
This appendix specifies the certificate entirely by integers and the actual
numbers of Section~\ref{sec:algebra}. All indices in the data are zero-based.
The factorized representation is more compact than printing the expanded
164 coefficients of the LDL-derived SOS, some of which have several hundred
digits. It changes neither those coefficients nor their weights.

\subsection{Phase supports}
Each row in Table~\ref{tab:phases} lists the four nonzero entries of column
$\ell$ of $F_k$. A pair $(j,e)$ means $(F_k)_{j\ell}=\zeta^e$; all unlisted
entries are zero. The row index $j$ refers exactly to \eqref{eq:q}.
\begin{longtable}{@{}rrl@{}}
\caption{Nonzero phase-support pairs $(j,e)$ in column $\ell$ of $F_k$.}\label{tab:phases}\\
\toprule $k$ & $\ell$ & $(j,e)$ \\ \midrule\endfirsthead
\toprule $k$ & $\ell$ & $(j,e)$ \\ \midrule\endhead
1 & 0 & $(1,0)\quad (5,2)\quad (9,19)\quad (45,1)$ \\
1 & 1 & $(22,0)\quad (42,7)\quad (62,2)\quad (74,1)$ \\
1 & 2 & $(26,0)\quad (38,5)\quad (58,6)\quad (78,7)$ \\
1 & 3 & $(30,0)\quad (34,11)\quad (66,1)\quad (70,2)$ \\
6 & 0 & $(1,0)\quad (5,12)\quad (9,14)\quad (45,6)$ \\
6 & 1 & $(22,0)\quad (42,2)\quad (62,12)\quad (74,6)$ \\
6 & 2 & $(26,0)\quad (38,10)\quad (58,16)\quad (78,2)$ \\
6 & 3 & $(30,0)\quad (34,6)\quad (66,6)\quad (70,12)$ \\
7 & 0 & $(2,0)\quad (6,14)\quad (18,13)\quad (54,7)$ \\
7 & 1 & $(10,0)\quad (14,17)\quad (46,7)\quad (50,14)$ \\
7 & 2 & $(31,0)\quad (43,5)\quad (71,14)\quad (75,7)$ \\
7 & 3 & $(35,0)\quad (39,11)\quad (67,18)\quad (79,5)$ \\
8 & 0 & $(3,0)\quad (7,16)\quad (27,12)\quad (63,8)$ \\
8 & 1 & $(11,0)\quad (23,16)\quad (51,16)\quad (55,8)$ \\
8 & 2 & $(15,0)\quad (19,0)\quad (47,8)\quad (59,16)$ \\
8 & 3 & $(40,0)\quad (44,8)\quad (76,8)\quad (80,16)$ \\
9 & 0 & $(4,0)\quad (8,18)\quad (36,11)\quad (72,9)$ \\
9 & 1 & $(12,0)\quad (32,15)\quad (52,18)\quad (64,9)$ \\
9 & 2 & $(16,0)\quad (28,17)\quad (48,6)\quad (68,15)$ \\
9 & 3 & $(20,0)\quad (24,19)\quad (56,9)\quad (60,18)$ \\
\bottomrule
\end{longtable}


\subsection{Small kernel factors}
Let $\langle n;d\rangle$ have the 24-coordinate meaning in
\eqref{eq:scalar}. The matrices $K_1,K_9$ have their last two rows equal
to $I_2$, and their first two rows are supplied below. The matrices
$K_7,K_8$ have their last three rows equal to $I_3$, and their first rows
are supplied below. Finally $K_6=I_4$. Every coordinate is printed, including
zeros. A line labelled $n[a{:}b]$ lists entries $n_a,\ldots,n_{b-1}$;
line breaks carry no further algebraic meaning.
\par\medskip\noindent\begin{minipage}{\textwidth}
\textbf{$K_{1}[0,0]$}
\begin{Verbatim}[fontsize=\fontsize{8}{9.4}\selectfont]
d = 456608800
n[00:06] = (-183905100, -201247900, -137909750, -49175550, 7673370, 4177738)
n[06:12] = (68247450, 20591550, 19432545, 7021015, -1042545, -337979)
n[12:18] = (-101720500, -34720100, -45724350, -9485230, 248530, 387514)
n[18:24] = (86040650, -6912450, 16514365, 4426015, -980665, -259591)
\end{Verbatim}
\end{minipage}

\par\medskip\noindent\begin{minipage}{\textwidth}
\textbf{$K_{1}[0,1]$}
\begin{Verbatim}[fontsize=\fontsize{8}{9.4}\selectfont]
d = 228304400
n[00:06] = (16933250, 28986350, -13620775, -1326085, 2374755, 832045)
n[06:12] = (4038550, -12864700, -14365490, -4193345, 333665, 128812)
n[12:18] = (-269277950, -33314850, -68067125, -29143135, 8258725, 3563555)
n[18:24] = (72093950, 2478800, 13797740, 3361605, -1793895, -718640)
\end{Verbatim}
\end{minipage}

\par\medskip\noindent\begin{minipage}{\textwidth}
\textbf{$K_{1}[1,0]$}
\begin{Verbatim}[fontsize=\fontsize{8}{9.4}\selectfont]
d = 228304400
n[00:06] = (22756100, -92877700, -4433400, -170460, -699060, 413880)
n[06:12] = (12403950, -3975850, 1556920, -1574790, 153420, -10786)
n[12:18] = (204063500, 58514700, 61269800, 29408400, -4842420, -2206444)
n[18:24] = (30784100, 12406500, 16181145, 7285315, -1182135, -501157)
\end{Verbatim}
\end{minipage}

\par\medskip\noindent\begin{minipage}{\textwidth}
\textbf{$K_{1}[1,1]$}
\begin{Verbatim}[fontsize=\fontsize{8}{9.4}\selectfont]
d = 456608800
n[00:06] = (-365083800, -7766400, -138101500, -60617840, 11548240, 4542004)
n[06:12] = (-17350650, -59665450, -34323185, -15510045, 2190215, 1095615)
n[12:18] = (-39403200, 150511800, 114618100, 37621400, -6459120, -4116452)
n[18:24] = (160719150, 1114950, 19091555, 12781095, -2436585, -420777)
\end{Verbatim}
\end{minipage}

\par\medskip\noindent\begin{minipage}{\textwidth}
\textbf{$K_{7}[0,0]$}
\begin{Verbatim}[fontsize=\fontsize{8}{9.4}\selectfont]
d = 400
n[00:06] = (-2400, -800, -700, -320, 60, 24)
n[06:12] = (500, 300, 205, 85, -15, -7)
n[12:18] = (-600, -600, -300, -160, 20, 16)
n[18:24] = (200, 200, 95, 55, -5, -5)
\end{Verbatim}
\end{minipage}

\par\medskip\noindent\begin{minipage}{\textwidth}
\textbf{$K_{7}[0,1]$}
\begin{Verbatim}[fontsize=\fontsize{8}{9.4}\selectfont]
d = 400
n[00:06] = (2900, 1300, 1000, 460, -80, -36)
n[06:12] = (-600, -300, -185, -85, 15, 7)
n[12:18] = (1500, 900, 500, 240, -40, -20)
n[18:24] = (-550, -250, -175, -95, 15, 7)
\end{Verbatim}
\end{minipage}

\par\medskip\noindent\begin{minipage}{\textwidth}
\textbf{$K_{7}[0,2]$}
\begin{Verbatim}[fontsize=\fontsize{8}{9.4}\selectfont]
d = 100
n[00:06] = (900, 400, 325, 135, -25, -11)
n[06:12] = (-225, -75, -60, -25, 5, 2)
n[12:18] = (0, 0, 0, 0, 0, 0)
n[18:24] = (0, 0, 0, 0, 0, 0)
\end{Verbatim}
\end{minipage}

\par\medskip\noindent\begin{minipage}{\textwidth}
\textbf{$K_{8}[0,0]$}
\begin{Verbatim}[fontsize=\fontsize{8}{9.4}\selectfont]
d = 200
n[00:06] = (250, -270, -90, -30, 4, 4)
n[06:12] = (0, 0, 0, 0, 0, 0)
n[12:18] = (0, 0, 0, 0, 0, 0)
n[18:24] = (-25, 25, 10, 0, 0, 0)
\end{Verbatim}
\end{minipage}

\par\medskip\noindent\begin{minipage}{\textwidth}
\textbf{$K_{8}[0,1]$}
\begin{Verbatim}[fontsize=\fontsize{8}{9.4}\selectfont]
d = 200
n[00:06] = (400, -20, 60, 40, -6, -2)
n[06:12] = (0, 0, 0, 0, 0, 0)
n[12:18] = (0, 0, 0, 0, 0, 0)
n[18:24] = (-175, 65, -5, -5, 2, 0)
\end{Verbatim}
\end{minipage}

\par\medskip\noindent\begin{minipage}{\textwidth}
\textbf{$K_{8}[0,2]$}
\begin{Verbatim}[fontsize=\fontsize{8}{9.4}\selectfont]
d = 1
n[00:06] = (0, 0, 0, 0, 0, 0)
n[06:12] = (0, 0, 0, 0, 0, 0)
n[12:18] = (0, 0, 0, 0, 0, 0)
n[18:24] = (0, 0, 0, 0, 0, 0)
\end{Verbatim}
\end{minipage}

\par\medskip\noindent\begin{minipage}{\textwidth}
\textbf{$K_{9}[0,0]$}
\begin{Verbatim}[fontsize=\fontsize{8}{9.4}\selectfont]
d = 456608800
n[00:06] = (240067500, 102352700, 66559350, 36328590, -3873570, -1331474)
n[06:12] = (15825550, 73550250, 35881365, 15315975, -1997665, -1224703)
n[12:18] = (7763100, -30074900, -25358950, 2861330, -203190, 265402)
n[18:24] = (-16752550, 3233950, -7512955, 731575, -99705, -79263)
\end{Verbatim}
\end{minipage}

\par\medskip\noindent\begin{minipage}{\textwidth}
\textbf{$K_{9}[0,1]$}
\begin{Verbatim}[fontsize=\fontsize{8}{9.4}\selectfont]
d = 456608800
n[00:06] = (7221100, -178266500, -131618450, -54794210, 7091910, 3325310)
n[06:12] = (21383350, -49300750, 5305245, -4555025, -287725, 121605)
n[12:18] = (-308864500, -90514500, -82301550, -36753110, 13455410, 6158450)
n[18:24] = (-100457950, -25109050, -30732945, -9485495, 3974165, 1679351)
\end{Verbatim}
\end{minipage}

\par\medskip\noindent\begin{minipage}{\textwidth}
\textbf{$K_{9}[1,0]$}
\begin{Verbatim}[fontsize=\fontsize{8}{9.4}\selectfont]
d = 456608800
n[00:06] = (-344881200, -96629000, -75496000, -34664220, 5696740, 2221680)
n[06:12] = (25865150, 32173150, 24601565, 7974285, -1597195, -730311)
n[12:18] = (-208560000, -156097000, -88408000, -48452940, 8045180, 3217592)
n[18:24] = (-75377850, 50503750, 27199135, 10667495, -1151845, -1119881)
\end{Verbatim}
\end{minipage}

\par\medskip\noindent\begin{minipage}{\textwidth}
\textbf{$K_{9}[1,1]$}
\begin{Verbatim}[fontsize=\fontsize{8}{9.4}\selectfont]
d = 456608800
n[00:06] = (-365083800, -7766400, -138101500, -60617840, 11548240, 4542004)
n[06:12] = (17350650, 59665450, 34323185, 15510045, -2190215, -1095615)
n[12:18] = (39403200, -150511800, -114618100, -37621400, 6459120, 4116452)
n[18:24] = (160719150, 1114950, 19091555, 12781095, -2436585, -420777)
\end{Verbatim}
\end{minipage}


\subsection{Hermitian parameters}
For each block, list the diagonal entries first, in increasing index, and
then the upper-triangular pairs $(a,b)$ in lexicographic order, giving
their real and imaginary parts in that order. The 42 real parameters
$h_0,\ldots,h_{41}$ have the assignments shown in Table~\ref{tab:assignments}.
An entry marked $D$ means $(H_k)_{aa}=h_m$; $R$ and $I$ mean respectively
$\Re(H_k)_{ab}=h_m$ and $\Im(H_k)_{ab}=h_m$ for $a<b$.
The lower triangle is the conjugate of the upper triangle.
\begin{longtable}{@{}rrrrcrrrrc@{}}
\caption{Parameter assignments. Each half-row is an independent record.}\label{tab:assignments}\\
\toprule $m$&$k$&$a$&$b$&type&$m$&$k$&$a$&$b$&type\\\midrule\endfirsthead
\toprule $m$&$k$&$a$&$b$&type&$m$&$k$&$a$&$b$&type\\\midrule\endhead
0 & 1 & 0 & 0 & D & 21 & 7 & 1 & 1 & D \\
1 & 1 & 1 & 1 & D & 22 & 7 & 2 & 2 & D \\
2 & 1 & 0 & 1 & R & 23 & 7 & 0 & 1 & R \\
3 & 1 & 0 & 1 & I & 24 & 7 & 0 & 1 & I \\
4 & 6 & 0 & 0 & D & 25 & 7 & 0 & 2 & R \\
5 & 6 & 1 & 1 & D & 26 & 7 & 0 & 2 & I \\
6 & 6 & 2 & 2 & D & 27 & 7 & 1 & 2 & R \\
7 & 6 & 3 & 3 & D & 28 & 7 & 1 & 2 & I \\
8 & 6 & 0 & 1 & R & 29 & 8 & 0 & 0 & D \\
9 & 6 & 0 & 1 & I & 30 & 8 & 1 & 1 & D \\
10 & 6 & 0 & 2 & R & 31 & 8 & 2 & 2 & D \\
11 & 6 & 0 & 2 & I & 32 & 8 & 0 & 1 & R \\
12 & 6 & 0 & 3 & R & 33 & 8 & 0 & 1 & I \\
13 & 6 & 0 & 3 & I & 34 & 8 & 0 & 2 & R \\
14 & 6 & 1 & 2 & R & 35 & 8 & 0 & 2 & I \\
15 & 6 & 1 & 2 & I & 36 & 8 & 1 & 2 & R \\
16 & 6 & 1 & 3 & R & 37 & 8 & 1 & 2 & I \\
17 & 6 & 1 & 3 & I & 38 & 9 & 0 & 0 & D \\
18 & 6 & 2 & 3 & R & 39 & 9 & 1 & 1 & D \\
19 & 6 & 2 & 3 & I & 40 & 9 & 0 & 1 & R \\
20 & 7 & 0 & 0 & D & 41 & 9 & 0 & 1 & I \\
\bottomrule
\end{longtable}


Every parameter has the shorter real encoding
\begin{equation}\label{eq:hencoding}
 h_m=\frac1{d_m}\sum_{\substack{a,c\in\{0,1\}\\ b\in\{0,1,2\}}}
 n^{(m)}_{a+2b+6c}s^ax^bu^c.
\end{equation}
The following records print $d_m$ and the twelve integers
$n^{(m)}_0,\ldots,n^{(m)}_{11}$ in order. For the purely rational entries
the same convention applies, with all omitted coordinates explicitly zero.
\par\medskip\noindent\begin{minipage}{\textwidth}
\textbf{$h_{0}$}
\begin{Verbatim}[fontsize=\fontsize{8}{9.4}\selectfont]
d = 1366860518953188827148918565372465257414896071591421359075903127985169665347544101683200
n[00] = 245412364478276689360459046688661315178543575230279807345743091655159290164819596824500
n[01] = -11394429692409728839452319816213821915646969738792970951239751826723221876667160201180
n[02] = 56058160874874101851967614584807804383309469172358857477518286579124223648777761130930
n[03] = 31661967076866684482740854152755631664517588918033470626290613896232261564362786744530
n[04] = -5226858198576092061421936943422072296776596065093349856130426689931098926096022811224
n[05] = -1719723616535715924221694811500546307738765675892967468032880510409167111695756291136
n[06] = -38605834107952402683397329884440923240923908037272537967439572587716408396244269952150
n[07] = 45855208658212317706561133915065813009197002395469800959083781641986039384509353814170
n[08] = 26901175805644363896060261641575276826569802757407685639064153041889781405942706728415
n[09] = 9683295717139454151924490072394208679804053932207163567867179044891007633054807136435
n[10] = -1233350356995512951723070765905916006705664807375737928974630570635826964551872298230
n[11] = -1049091332021440928986200466351893255994130298396244054580933671270589306231841979944
\end{Verbatim}
\end{minipage}

\par\medskip\noindent\begin{minipage}{\textwidth}
\textbf{$h_{1}$}
\begin{Verbatim}[fontsize=\fontsize{8}{9.4}\selectfont]
d = 2147483648
n = (8225705, 0, 0, 0, 0, 0, 0, 0, 0, 0, 0, 0)
\end{Verbatim}
\end{minipage}

\par\medskip\noindent\begin{minipage}{\textwidth}
\textbf{$h_{2}$}
\begin{Verbatim}[fontsize=\fontsize{8}{9.4}\selectfont]
d = 2733721037906377654297837130744930514829792143182842718151806255970339330695088203366400
n[00] = -38256020567417341341675935462905264097332864255380878686040986027665953187429911225200
n[01] = -24606636317723278746244325175715521490182605734409457281331215790702412580705344721360
n[02] = -60554357265422737177103669874149553352329903041547880640517248110054950862557931525080
n[03] = -28106984351329402664736409169512793971052066007698731331993030607179170521905989508280
n[04] = 4413427707546670504649328514582767041279239493047599573781384528512538456551163049412
n[05] = 1907539673572048163227221155676084906632126936325139347094329821457397427179335992312
n[06] = -44212618359121564358811850356015848982352863537135950066021767523854188050398232943100
n[07] = -12743989248689301359360486584387309602131602520587683467925602763179741358021880281460
n[08] = -31940318262999174876467224492355346155416458657343664985312116383261812196283222354550
n[09] = -13431330681908117164275203567081834201251756494773580281089428620499707082467663470470
n[10] = 2162100534425889937854093157603641648780116623498056984048582011195677377965316246685
n[11] = 1001955597509242099796060582059569919429599821161873314447451312427347954241591241525
\end{Verbatim}
\end{minipage}

\par\medskip\noindent\begin{minipage}{\textwidth}
\textbf{$h_{3}$}
\begin{Verbatim}[fontsize=\fontsize{8}{9.4}\selectfont]
d = 2733721037906377654297837130744930514829792143182842718151806255970339330695088203366400
n[00] = 419844232216811423813115793811669844798897141823675698375824496757701896246588308124800
n[01] = 30602732617905862080206452147102474994091296844914153295507104501815697258355473849760
n[02] = 79831192706184075974801872078511085259433518473893702190482949571491310598596848890100
n[03] = 40611801169467291914847744761090599253997081241859863428698150308632666005519216736900
n[04] = -7455284672914820700599319752054110883475898464840922868610504796319836106166956731772
n[05] = -2674345528839111162451452009992288746985064751927469096307125585390634156051123804328
n[06] = -24182000267372792064399782931127044108098835832959990080771322226080304048968841417700
n[07] = 45463681556841011658977730385142842743614474390280970761882119898378183899242043266580
n[08] = 4603572073360186226454672229629754366557429993684504456330604386161542631393146401960
n[09] = -193074976580091783317615567117509413882702734560802347658649687532051644390634221300
n[10] = -35658237387427591033925126504862628441273385735054396555433834115733564747658665709
n[11] = -265384743841121472141337163572908168796665491182932452130957252744629249741520905405
\end{Verbatim}
\end{minipage}

\par\medskip\noindent\begin{minipage}{\textwidth}
\textbf{$h_{4}$}
\begin{Verbatim}[fontsize=\fontsize{8}{9.4}\selectfont]
d = 4294967296
n = (501845005, 0, 0, 0, 0, 0, 0, 0, 0, 0, 0, 0)
\end{Verbatim}
\end{minipage}

\par\medskip\noindent\begin{minipage}{\textwidth}
\textbf{$h_{5}$}
\begin{Verbatim}[fontsize=\fontsize{8}{9.4}\selectfont]
d = 341715129738297206787229641343116314353724017897855339768975781996292416336886025420800
n[00] = -93212176226471784315346243448540559929781573142736056249018851934681535418905441421400
n[01] = -37440551519353710862408995506089595206083555776081770316545176523686031181754268862480
n[02] = -25344201712487537381589900793295065046574769329082571531232772907258850493444635493695
n[03] = -11992972421091229333263957115531600660200011667902366227216530794394681035167130734275
n[04] = 2090932940489596663134590636357937150161814468430108266393418332051713913225110915826
n[05] = 894032334677204332150915941475923370477422441172010071712996389426285014400070559512
n[06] = -20726383164457248371023186826754607060452955428765565093140707463428856793322457763700
n[07] = -10631034114272716935040580232479770182037436837895224175128912241958827012917469703410
n[08] = -4350803428849928180278364478869773820568238972306835999079431395417517862018113460285
n[09] = -1848570118441532880588348840435462602588315611712963648393759866728098538632274085035
n[10] = 438375764953295351443099316143887355415037405665708597169441270567905117620939555240
n[11] = 161725756314982541383946137710294653171640357058486212237190130971295558619718442369
\end{Verbatim}
\end{minipage}

\par\medskip\noindent\begin{minipage}{\textwidth}
\textbf{$h_{6}$}
\begin{Verbatim}[fontsize=\fontsize{8}{9.4}\selectfont]
d = 683430259476594413574459282686232628707448035795710679537951563992584832673772050841600
n[00] = -82813245166154708832530792263294858337154967914530818407238643214912365249392405920200
n[01] = -112727247288116951468484859652223835094220474241478008693022711737892627976435294631840
n[02] = -52163633280533279436566246868021086298825155525945100442418776995851283891938974425900
n[03] = -21123027888442893218012284526636632207824247009397506214731817839508550279748993786620
n[04] = 3389069931730817865653006893327215321063274459578432605717808341633093995057785072798
n[05] = 1955371905279809757547724663732558062461821980915374371023134086028305955355463478514
n[06] = -81380726078805152344178444828925233708139081557877321082959929039593559103860004898350
n[07] = -5637510202171919454390278658442929944542413341409305940996761993169988807750107641550
n[08] = -7769186134673623995441339772116722697915177268368749649817691352871339985133353177235
n[09] = -4693270051363230320989851053850162150334390163749034311331160295773416100135786128265
n[10] = 1169058040498867847511648937981339172690888030080882221980300999461714119039492277101
n[11] = 241947511099134816754551773183040713160661976060041431877147654019255366879786166501
\end{Verbatim}
\end{minipage}

\par\medskip\noindent\begin{minipage}{\textwidth}
\textbf{$h_{7}$}
\begin{Verbatim}[fontsize=\fontsize{8}{9.4}\selectfont]
d = 683430259476594413574459282686232628707448035795710679537951563992584832673772050841600
n[00] = 395280926859635089268287986122953562473577534344063262628042631127107059845859692244000
n[01] = 202045504877653826245704078244601253809701559157461368006861870684162863481805673870960
n[02] = 173185154483668746276181970756186079947428477596572793950067573631171625861338026066750
n[03] = 76948919657316977313367689998831652795894243016658504639882861733206992117144340039370
n[04] = -13605535734670015213895172119244421558827243279858456462998045889027072666316395979092
n[05] = -6067854269979903276078973862181132061117673023085767664554193459298329711890310045796
n[06] = 152003786003313447808793420385613513315762402568179271373568283548903999748792021314300
n[07] = 41080346192109271184710294659061440275348863412584052603094052612152746487579243195300
n[08] = 53742953141309790077928198076107108137353139255578571396508577600613548805649853610700
n[09] = 21857343706088014070490496918952873721913879266235851169180725805617860633164647504660
n[10] = -3919508062073923341517111566683315382619502949671728208536392467092888700431970178443
n[11] = -1703758227107723189941572893055163742852063087979235472354013635557531142341947506755
\end{Verbatim}
\end{minipage}

\par\medskip\noindent\begin{minipage}{\textwidth}
\textbf{$h_{8}$}
\begin{Verbatim}[fontsize=\fontsize{8}{9.4}\selectfont]
d = 4294967296
n = (-39622637, 0, 0, 0, 0, 0, 0, 0, 0, 0, 0, 0)
\end{Verbatim}
\end{minipage}

\par\medskip\noindent\begin{minipage}{\textwidth}
\textbf{$h_{9}$}
\begin{Verbatim}[fontsize=\fontsize{8}{9.4}\selectfont]
d = 2147483648
n = (-17650465, 0, 0, 0, 0, 0, 0, 0, 0, 0, 0, 0)
\end{Verbatim}
\end{minipage}

\par\medskip\noindent\begin{minipage}{\textwidth}
\textbf{$h_{10}$}
\begin{Verbatim}[fontsize=\fontsize{8}{9.4}\selectfont]
d = 2793135472143228293694245096652800
n[00] = 47835841622898479865456095142400
n[01] = 5493272592418579007118902408760
n[02] = 4066215439620936259187657286770
n[03] = 3831266685103962067343222653570
n[04] = -744356613373203440029482562062
n[05] = -323849719680936246223857542582
n[06] = -22478579129120105799812171276350
n[07] = 4615916590280372440756468702950
n[08] = -556272474815457729689662127855
n[09] = -1236347750335889564584875593545
n[10] = 327915155095382705206506573515
n[11] = 25985190790067278933108692577
\end{Verbatim}
\end{minipage}

\par\medskip\noindent\begin{minipage}{\textwidth}
\textbf{$h_{11}$}
\begin{Verbatim}[fontsize=\fontsize{8}{9.4}\selectfont]
d = 2147483648
n = (24295977, 0, 0, 0, 0, 0, 0, 0, 0, 0, 0, 0)
\end{Verbatim}
\end{minipage}

\par\medskip\noindent\begin{minipage}{\textwidth}
\textbf{$h_{12}$}
\begin{Verbatim}[fontsize=\fontsize{8}{9.4}\selectfont]
d = 2147483648
n = (-17690197, 0, 0, 0, 0, 0, 0, 0, 0, 0, 0, 0)
\end{Verbatim}
\end{minipage}

\par\medskip\noindent\begin{minipage}{\textwidth}
\textbf{$h_{13}$}
\begin{Verbatim}[fontsize=\fontsize{8}{9.4}\selectfont]
d = 4294967296
n = (-48696935, 0, 0, 0, 0, 0, 0, 0, 0, 0, 0, 0)
\end{Verbatim}
\end{minipage}

\par\medskip\noindent\begin{minipage}{\textwidth}
\textbf{$h_{14}$}
\begin{Verbatim}[fontsize=\fontsize{8}{9.4}\selectfont]
d = 1366860518953188827148918565372465257414896071591421359075903127985169665347544101683200
n[00] = 491821951719985230069443879781708176212733420554960115133841345419262842448534338427300
n[01] = 391966458315792933197496766401906700863722949365935764349969416363424225073274916892420
n[02] = 301256234711809971710813505121445251107358423477840672778260189363626312354935474133270
n[03] = 129173309607949308513722236976252976776571091748984855606724835628698343181097844670090
n[04] = -20971663187286138533175998314884294022055526213630269807080674073661533327391857753794
n[05] = -10386493618585598196666807950195989279621189514316448550886345627209297704027982731346
n[06] = 225888667140019892593042155692138293887962013307566903119903630530446436541845138370900
n[07] = 37908700072007305964663447672894843756603477579083858084612152763167042039559346395980
n[08] = 55550270252982232561800598649185918056200511357092145271711725045494868589691015560530
n[09] = 25070968389394611401262962913162787624591450748132234033179543400573364568071836239530
n[10] = -4898198755646738653034172716969288674509453750360464825232230159001885796701824059039
n[11] = -1840526206698916110324133784395320281915074966737362116697277295410924558339124573913
\end{Verbatim}
\end{minipage}

\par\medskip\noindent\begin{minipage}{\textwidth}
\textbf{$h_{15}$}
\begin{Verbatim}[fontsize=\fontsize{8}{9.4}\selectfont]
d = 1366860518953188827148918565372465257414896071591421359075903127985169665347544101683200
n[00] = -701169881350883339136082240743956064022576076654121133829045278183804642183475586615200
n[01] = -305981895736927934280400325350050764040076937791788385416993795961842352203302128545040
n[02] = -322373980083975407649029745562292559020006643262354285101265367243473986946607628524940
n[03] = -145416338438313811820022023366501584844735710021704450895359533131805745844932046532460
n[04] = 23834704128278657859159462713747361172142627303509887718205894277112014118672437626838
n[05] = 10943123393189751123205138157425143401219284167683889225012605796971175740850000541582
n[06] = -176698055856078159586379366804295204309672878988334705582068303766714329462069869416250
n[07] = -89845631121011001067129858260577779179184449122791608454148921718139758366437747165350
n[08] = -88979504489797271035899026375392859741256986121016026690198645668090324620392038975460
n[09] = -37749366648935951105950978381035589335958050493143073275220756288556885716336916329050
n[10] = 6462988556525557135082145361250342666575314533312194968350987501073163694599824540649
n[11] = 3080182308820291525720808492637703706105361340643801865062432764188636510535112892539
\end{Verbatim}
\end{minipage}

\par\medskip\noindent\begin{minipage}{\textwidth}
\textbf{$h_{16}$}
\begin{Verbatim}[fontsize=\fontsize{8}{9.4}\selectfont]
d = 1366860518953188827148918565372465257414896071591421359075903127985169665347544101683200
n[00] = -293648548014998808178401895229195188994139374251549960306003719060731006360240909121850
n[01] = -64440782780400855492497328077648243319858115870263257675019157911420090162664870713990
n[02] = -76400626051462657133363682218541365934785412993732449097691507375093979788022878264695
n[03] = -34997030896547601817560935007135811351024045639518675432216230858476424686614518325705
n[04] = 6833876329074046562490838720378342154030942108031880924098559086509730518492789668433
n[05] = 2776286273956831781476907972703293691098915074610924755143590359469417260184958456235
n[06] = -34719511464791184178908990098975510729971531290033933150606959143331581482894774804650
n[07] = -23718425791162730557130778253663049973389213523776108056725968044781445924469524389360
n[08] = -8824663749148716559693550171113744627761146661750150678403587285232987721685757244505
n[09] = -2221391246295701535939451368526492102010124543742593688520109904225173894158167302960
n[10] = 436418029367253624903500102374836330250440155972559795625713506180864923291584007586
n[11] = 286535796221470885157363030910618814751463590437079463637759452609837078163652559543
\end{Verbatim}
\end{minipage}

\par\medskip\noindent\begin{minipage}{\textwidth}
\textbf{$h_{17}$}
\begin{Verbatim}[fontsize=\fontsize{8}{9.4}\selectfont]
d = 1366860518953188827148918565372465257414896071591421359075903127985169665347544101683200
n[00] = 704487079928962447252273807350476610302335617985672392059171714294392208670796192403650
n[01] = 215280937775993329145406025873464891268060881126252658291206624383265586771850575132830
n[02] = 258049533463973200508104340008053711114653720193372322850712078595150146705730201012305
n[03] = 121737155934141381553582116030735994968204128690047416801834169304041104917155808943975
n[04] = -21152292729336375714064715190591200371366900263312796120377409303603674197739672585131
n[05] = -8717995187861924567890250886078635944796785881921510109055011462209230194554950462545
n[06] = 93233095259324286549192837682886936467158221385026085967912511099918618188626849990750
n[07] = 95018019867110961974560408422078036226296011194495341823299450379975383897626725598760
n[08] = 70677462892574006265995604971217947745922753395038946690322511180410096136229280275135
n[09] = 31114694028205426076220744387996342515732501853468311032453744891743306601841845041810
n[10] = -4997114845233648960711746350821699917036222782105640770097246236046739685283823205752
n[11] = -2492668135280212470632317338653195315259108075197621056651502284415109055378495291713
\end{Verbatim}
\end{minipage}

\par\medskip\noindent\begin{minipage}{\textwidth}
\textbf{$h_{18}$}
\begin{Verbatim}[fontsize=\fontsize{8}{9.4}\selectfont]
d = 1366860518953188827148918565372465257414896071591421359075903127985169665347544101683200
n[00] = 289837895602342935367166540911987810456706027552505472752243882164674783511641391554450
n[01] = 66311959030153299885848170978575168388516926613721588300611860093195157542004365684510
n[02] = 78261923961813475654599866165394096475508904068442078745277496495017818547300628678135
n[03] = 34940874102476541928788225558631821071024793095602389107651857549826137714130148820365
n[04] = -6796859062051714835640144604915000233957349709188068250066867998915772388625779560697
n[05] = -2848835304525121936466544613904776465765568277434587890679837943942673666224659959319
n[06] = 38673379646690193170074167540314996712760192381710603836575273353708121745635186371300
n[07] = 21906303454447974971452232233129905338259244477607381372893421180862339615809036237630
n[08] = 8017814461449459370302171746063288561567462454113207297923947437342603742594121299765
n[09] = 2378178214838578429203396434447919509874818297275400009343415142700955470525688198290
n[10] = -471160078230348891121105283533177014129862517045967536909147693659556052889036022050
n[11] = -256279679297541068734814130344307091878015975138899726326201498290608482512629515657
\end{Verbatim}
\end{minipage}

\par\medskip\noindent\begin{minipage}{\textwidth}
\textbf{$h_{19}$}
\begin{Verbatim}[fontsize=\fontsize{8}{9.4}\selectfont]
d = 1366860518953188827148918565372465257414896071591421359075903127985169665347544101683200
n[00] = 507153726821380688223057687137540022360364471407773717535690271416863228577195683920150
n[01] = 283046041768009977951679207002782068848650946086926768100153096064910857443325137836130
n[02] = 258092712763633830624398433235594412777077187031456127790728201262373613703001735797655
n[03] = 115470538676917693114094925155392270065188777606512408876293727523922131688687710491545
n[04] = -19537484576626436751502415032259450588995179087244827008678620231234355210108735266701
n[05] = -8926214332351986306240355680816690531745488846137868996336763127986115780180618187191
n[06] = 174794051483108262589507981551482777704737012125033756107497935090215255438122440970450
n[07] = 63927666587664211760086815606017447080809097385734175748745241289107963405230015048140
n[08] = 69761869221950369037138543668176297365308109756794205265233275530114907397453894492300
n[09] = 33166361756681887316258396828933015756740127882392914669805684579351166381626075082295
n[10] = -5621381934454402294718865621121137306422317204026373785618083807290002682191642329133
n[11] = -2348585847556249265506878394444337405721048507247443281453563730116445724673684702950
\end{Verbatim}
\end{minipage}

\par\medskip\noindent\begin{minipage}{\textwidth}
\textbf{$h_{20}$}
\begin{Verbatim}[fontsize=\fontsize{8}{9.4}\selectfont]
d = 1366860518953188827148918565372465257414896071591421359075903127985169665347544101683200
n[00] = -235982689539709975972991465445007491550745230253188468677834392359208605474052005642700
n[01] = -164552547714143411506728270696638716152862363769146482162759461935357887521215447929100
n[02] = -115026484076358557141398516482095289974940960111404727971975811276212185765203277500870
n[03] = -45406665759829249727064809883614224725355289654813822575992874948783748147570751747870
n[04] = 8403057552118811608509887009003120138005627004253164762839887460970105502899456977770
n[05] = 4455395253145158237681350227826603320253844388103827322631047603106423215184838858738
n[06] = -112555937775662311162696769785031672069306921498997614642240548750755974013411441501550
n[07] = -12345442507332585503562864811002971548844180255758377921632756344033911115118078036490
n[08] = -54962554299837920387718071638899590782703880258150495883721863869035353137891049226805
n[09] = -22502308497553752662064734161914421450988809732737277833368503669750173127993161453495
n[10] = 3577349163265773503047826141361266901375947402153818478014443139260590576132012224691
n[11] = 1640917197794259429716225825129027098875393705512955962813041207391469670287038035345
\end{Verbatim}
\end{minipage}

\par\medskip\noindent\begin{minipage}{\textwidth}
\textbf{$h_{21}$}
\begin{Verbatim}[fontsize=\fontsize{8}{9.4}\selectfont]
d = 8379406416429684881082735289958400
n[00] = 296050830919368131257441869618650
n[01] = -151732213946955710742508975434050
n[02] = -37552369609800877309515350858775
n[03] = -19239973838604527732732026729965
n[04] = 1771162365165672556769707065405
n[05] = 1802110731117904800417567881591
n[06] = -107151074367379930925675642768850
n[07] = 57518642834688666046406695661800
n[08] = 10344738626542428095901441853590
n[09] = 4845764606519362846077015943335
n[10] = -460005211828225354299233114600
n[11] = -476204889207908603032733650067
\end{Verbatim}
\end{minipage}

\par\medskip\noindent\begin{minipage}{\textwidth}
\textbf{$h_{22}$}
\begin{Verbatim}[fontsize=\fontsize{8}{9.4}\selectfont]
d = 4294967296
n = (19694267, 0, 0, 0, 0, 0, 0, 0, 0, 0, 0, 0)
\end{Verbatim}
\end{minipage}

\par\medskip\noindent\begin{minipage}{\textwidth}
\textbf{$h_{23}$}
\begin{Verbatim}[fontsize=\fontsize{8}{9.4}\selectfont]
d = 2733721037906377654297837130744930514829792143182842718151806255970339330695088203366400
n[00] = 243724321419008311235385953077050887168464483820632996153070523023829526776764263685100
n[01] = 199146037199101695294576522333348032267329654498498501534842764254632528988952757336300
n[02] = 99552237782962777945734446575668234192316147426865165066142025736453823983509995947400
n[03] = 41834960618374495916923448606135598267214692703656650871110670698147980108871831462700
n[04] = -7521987311761556960596494538087691612338364731092313670582138591537998935526070960760
n[05] = -3738239736074585498384047768421480280712283333802187493597341120949346362748840653136
n[06] = 171819277401554044884462553306513336223771118192081862103303679190458550351539572469550
n[07] = 28382680992492032227575868910420724874864895485832845639664530387481429244630727082770
n[08] = 22718160367973672521699543570581690080529234476881013406836945389499578569514041890505
n[09] = 11408442049685906353588991812702818440997204486528343055042288115232218634879338057105
n[10] = -2472727747014319016917806964329164875742922562499039620066170921204923787024848410216
n[11] = -783120868471591524536529034452981675286819806984698515943462603720103226019315517636
\end{Verbatim}
\end{minipage}

\par\medskip\noindent\begin{minipage}{\textwidth}
\textbf{$h_{24}$}
\begin{Verbatim}[fontsize=\fontsize{8}{9.4}\selectfont]
d = 2733721037906377654297837130744930514829792143182842718151806255970339330695088203366400
n[00] = -659188627283393470512733592499641906893164294771735430705696036614546822723144701270300
n[01] = -259816354186404115859253078296752440755848093868815214321588302892663690875074538393260
n[02] = -191920939487370320371489209019022172362102112967996585430210831768122114773561538100820
n[03] = -91802805894742434779527963429813807393066687458554366569460166694873153403719189608840
n[04] = 15853136588607838984393618921762867076111076420295619602206316508024707173268052486152
n[05] = 6623424020982660277730332140792015589887488810705752924452205083895036950582903034424
n[06] = -145950948389173562104132836109465659836888028605918061856359150113004781191759697996150
n[07] = -81151669743461939948455863961458246302868479179168668437820775234253671339549240852210
n[08] = -59630391970241955172230381490731748445658877169359429225009077650236727470371426246465
n[09] = -25180777496992795375231065818777820357882199977503636714867896861643907090539943585405
n[10] = 4549894098731050180205023458057994500967369588491081787207102240086730571663703467830
n[11] = 2179125837240238013749139260713042491567159230439507073333272015833025166955934105010
\end{Verbatim}
\end{minipage}

\par\medskip\noindent\begin{minipage}{\textwidth}
\textbf{$h_{25}$}
\begin{Verbatim}[fontsize=\fontsize{8}{9.4}\selectfont]
d = 2733721037906377654297837130744930514829792143182842718151806255970339330695088203366400
n[00] = -250553033582924209256426150408920727144255746910931137716500667781517906788371960996300
n[01] = -97055588281364852536194070669251333826434367052840461602312372905146990344534161130540
n[02] = -17337300432961584193387031954335202013680017342179711669289307701612185497251890210030
n[03] = -10505683296377681936667029911080871395620173770369658306833325674104549980695497813730
n[04] = 2844762035474905405786361689229583142869223892718926593102662634174708324489136589128
n[05] = 1063172078339764659609138723709120920926405368353453786134236967761129962828389432236
n[06] = -29904819888617837280643272760277470710587632606086009957824265433726143546325984113550
n[07] = -18469203855839912357817877087221630210804114919387908633971666565217113762094808508690
n[08] = -7757097367572246748080024992388030286025761089684671175223220740516952779436030239230
n[09] = -3381842134154483737490474373055556391581940447114444565278417382369610362436524478160
n[10] = 577996338716248034560052937993898805727633763740820235180786510220097655764940592215
n[11] = 283582895246441484159006166620564107181528025983396965581205966817865967248801328455
\end{Verbatim}
\end{minipage}

\par\medskip\noindent\begin{minipage}{\textwidth}
\textbf{$h_{26}$}
\begin{Verbatim}[fontsize=\fontsize{8}{9.4}\selectfont]
d = 2733721037906377654297837130744930514829792143182842718151806255970339330695088203366400
n[00] = -562830155022568897629358425563413156252792496788415731847626867032696421595723257980100
n[01] = -359032342063169199247281779889368672817166625430337352237955892000102703424025166342100
n[02] = -236543800968016981470324478702398801816351639885427629898783024877299966028111582932950
n[03] = -102564360167434091233976295383880572925689476951728479201729976983696404315125620328090
n[04] = 18191253781606707323580234753220770489511999717885174839086578303344937588652328117020
n[05] = 8808866076212798450921021798204154780332499483039180356351205710651551107607143721576
n[06] = -250574890939257586462576380233491278546891903941910724128957769106385289549847122229850
n[07] = -62853693558685941733355135881914871359766790054827194086270554223273514049683018252990
n[08] = -56300583958326366899116907744908144081780983129717066022495356213888842138904176950110
n[09] = -27452034649086269009397842661534618903825940686188064704134986125021360718065460810300
n[10] = 4927128202460291736655595430187529563121710910130799656874272994431631251612337015725
n[11] = 1939040937199528593396101691355621151103612660494924596578833331471793275933020478657
\end{Verbatim}
\end{minipage}

\par\medskip\noindent\begin{minipage}{\textwidth}
\textbf{$h_{27}$}
\begin{Verbatim}[fontsize=\fontsize{8}{9.4}\selectfont]
d = 2733721037906377654297837130744930514829792143182842718151806255970339330695088203366400
n[00] = -2169085024488417649901865659703565209907645304173251361583994403858819489934943850717100
n[01] = -1021473009998123196012974458889945618030638825313803213358792758599000732841411000982300
n[02] = -748165241270811700358680859795605328036721354910894823732917904454286687508744953254750
n[03] = -346987449762530292497210588365813313531968224721285089602529211575120694714459218687270
n[04] = 60945677283331825058186750224962731484440651461644897793064068908920851909605351657290
n[05] = 26753113812511773997855959862625635363734638561330330179303522774044011308110251592042
n[06] = -565268468251931844765221740909132618519472523320346385351449945766214220021749067832650
n[07] = -270880509254750400217510676139370956690757901860328414147873665186986197568134415187750
n[08] = -207182615826336156304479034060747723787412768434695140340268965947679130655925597994865
n[09] = -97808533790309628480408859627749453505348095431899186439604349599131887212594072074295
n[10] = 16985157364679875615315982744048857932392096574002929460657775014307202363761195376649
n[11] = 7329612373738830774395503298504145590971351169958801550526933162907716820352127681771
\end{Verbatim}
\end{minipage}

\par\medskip\noindent\begin{minipage}{\textwidth}
\textbf{$h_{28}$}
\begin{Verbatim}[fontsize=\fontsize{8}{9.4}\selectfont]
d = 2733721037906377654297837130744930514829792143182842718151806255970339330695088203366400
n[00] = 703674474812469337837034920473801212020224723393024733518630695257242645689760171878900
n[01] = 377431124057514078684655770621443185490614934522567563087047469879112364383582125677300
n[02] = 283307000000753038114030280079713994997088276072112100306339922956060879049173632010530
n[03] = 124576817962959132327063908833614890027299592290158032118791728332032038593751105989090
n[04] = -21501592834756657098796414920168205297962140897051539133476201493268928826721979666630
n[05] = -9994372573992388509312464552850544255506299933350815391982605176907736103618906434606
n[06] = 212342936250949383107560442468479860534292216435332913564712781331782940842005398686350
n[07] = 72780105256824095732142185833875073186281030608148404495630471835959458090443305486250
n[08] = 74989220926855767446876191295393862146408750233179873553154005292538967670391963756855
n[09] = 34190324350208904394272436161120939015139882025364810625564396112599032935263700791845
n[10] = -6042518418324536064268547612114308782062374041819670228394847975747865604314060755127
n[11] = -2551673843885838760627063095913947032159703430800600925808470636724398511839040153297
\end{Verbatim}
\end{minipage}

\par\medskip\noindent\begin{minipage}{\textwidth}
\textbf{$h_{29}$}
\begin{Verbatim}[fontsize=\fontsize{8}{9.4}\selectfont]
d = 4294967296
n = (19694267, 0, 0, 0, 0, 0, 0, 0, 0, 0, 0, 0)
\end{Verbatim}
\end{minipage}

\par\medskip\noindent\begin{minipage}{\textwidth}
\textbf{$h_{30}$}
\begin{Verbatim}[fontsize=\fontsize{8}{9.4}\selectfont]
d = 8379406416429684881082735289958400
n[00] = 296050830919368131257441869618650
n[01] = -151732213946955710742508975434050
n[02] = -37552369609800877309515350858775
n[03] = -19239973838604527732732026729965
n[04] = 1771162365165672556769707065405
n[05] = 1802110731117904800417567881591
n[06] = -107151074367379930925675642768850
n[07] = 57518642834688666046406695661800
n[08] = 10344738626542428095901441853590
n[09] = 4845764606519362846077015943335
n[10] = -460005211828225354299233114600
n[11] = -476204889207908603032733650067
\end{Verbatim}
\end{minipage}

\par\medskip\noindent\begin{minipage}{\textwidth}
\textbf{$h_{31}$}
\begin{Verbatim}[fontsize=\fontsize{8}{9.4}\selectfont]
d = 1366860518953188827148918565372465257414896071591421359075903127985169665347544101683200
n[00] = -235982689539709975972991465445007491550745230253188468677834392359208605474052005642700
n[01] = -164552547714143411506728270696638716152862363769146482162759461935357887521215447929100
n[02] = -115026484076358557141398516482095289974940960111404727971975811276212185765203277500870
n[03] = -45406665759829249727064809883614224725355289654813822575992874948783748147570751747870
n[04] = 8403057552118811608509887009003120138005627004253164762839887460970105502899456977770
n[05] = 4455395253145158237681350227826603320253844388103827322631047603106423215184838858738
n[06] = -112555937775662311162696769785031672069306921498997614642240548750755974013411441501550
n[07] = -12345442507332585503562864811002971548844180255758377921632756344033911115118078036490
n[08] = -54962554299837920387718071638899590782703880258150495883721863869035353137891049226805
n[09] = -22502308497553752662064734161914421450988809732737277833368503669750173127993161453495
n[10] = 3577349163265773503047826141361266901375947402153818478014443139260590576132012224691
n[11] = 1640917197794259429716225825129027098875393705512955962813041207391469670287038035345
\end{Verbatim}
\end{minipage}

\par\medskip\noindent\begin{minipage}{\textwidth}
\textbf{$h_{32}$}
\begin{Verbatim}[fontsize=\fontsize{8}{9.4}\selectfont]
d = 2733721037906377654297837130744930514829792143182842718151806255970339330695088203366400
n[00] = 1794502157397930348560270059462905259167363575646733728269157363848421249917671093424100
n[01] = 878274669574103108138292773339904194354228569093379935437513329506040033604666060335300
n[02] = 677583601588103806081897418199152829447049738321816301877778609191002716191403700689170
n[03] = 308430096879493672906520073288222719414629451851607019222258308315662073245592347677810
n[04] = -53669356837245444612337238145361379813516279667478693071352564345219860037440169519450
n[05] = -23939804551495947596275762026455377683010465970973287671272821335332595051985148584746
n[06] = 504381858613811638587178793250667426127633091891960563167638706502701285081183180493050
n[07] = 220477544750999477159389050563835207670656909871654676072427612275794312522462226945550
n[08] = 163050710111655736458548717571348623776857673754312661039562482295957622625704603263095
n[09] = 76547138296470857361151708307885097600174839128367506668734900598795189994832738930565
n[10] = -13557456620556791829829995589376826276426127087365390848104140925261681238681052911483
n[11] = -5815567309518269173553618836606318403457991987577815219179286358755136053908807747553
\end{Verbatim}
\end{minipage}

\par\medskip\noindent\begin{minipage}{\textwidth}
\textbf{$h_{33}$}
\begin{Verbatim}[fontsize=\fontsize{8}{9.4}\selectfont]
d = 2733721037906377654297837130744930514829792143182842718151806255970339330695088203366400
n[00] = -1447377610144983279640008229442428054362820323207644471453350480203402040899149548980100
n[01] = -575024734890108544356358486022332511202100304452899505077732981570744804920494175492500
n[02] = -519641865078121620134714589149652312809651522820244366516728041868174198030026946417150
n[03] = -223264869216020496674486753907947282237242362727404396873085166147558564371926667499910
n[04] = 39690453600332698295012289042112000869113329017802366161371682771505461448119228846790
n[05] = 17948584686581877698784629136426296001606948822989979332158737113404139232128405777118
n[06] = -373202138208572373539841640065380844238635427343580104726637268856489853377170794996250
n[07] = -167954770763673880808091708847801211829832388995646383572655937614585096709299194580550
n[08] = -141291763281491256052898886039428384684103996199228223041023176251010296114054598596785
n[09] = -58487724999746415037783520816653290077341066323238496504864086170144820509270657977415
n[10] = 10291124334692733838864487167160018830106700043210654356363273073375077641180355674709
n[11] = 4912479391233358337558235999598814149231761334348735825528361757076805411756993532371
\end{Verbatim}
\end{minipage}

\par\medskip\noindent\begin{minipage}{\textwidth}
\textbf{$h_{34}$}
\begin{Verbatim}[fontsize=\fontsize{8}{9.4}\selectfont]
d = 2733721037906377654297837130744930514829792143182842718151806255970339330695088203366400
n[00] = 429017948184463648747915493089605304261739526404502553815027978239312417152152624988200
n[01] = 112074872823758624130385490516703931741210571441032371466259042888054285721324536238000
n[02] = 96916848721605194888815497967859982397043190060301480878720810317511921741560516470350
n[03] = 45828014219318055276853244258462463344633272396053234675070915799583392724055816966230
n[04] = -8617217234771045029553866320986158135310528898638224737572343755073113544880202424390
n[05] = -3343552333822704945842345124302513598112512005540401152227251320328060366265769716106
n[06] = 78561652609274082865156362396250951349049448310789096004723917447875046871734982507850
n[07] = 71194196415484124316353828755206935253389870234980997911249896947064691461174566276850
n[08] = 24574222430016640585314834379275038112757184410850742291867220528207536737168754327855
n[09] = 9838558151404444956596523748613599143393804053474664906298423940743007850383553488385
n[10] = -1903209629753064159027818678955039833816394071265687479780639074275510863135475491672
n[11] = -1012814835900131950717158115635257333236125904497332211607419157680241984626926492326
\end{Verbatim}
\end{minipage}

\par\medskip\noindent\begin{minipage}{\textwidth}
\textbf{$h_{35}$}
\begin{Verbatim}[fontsize=\fontsize{8}{9.4}\selectfont]
d = 2733721037906377654297837130744930514829792143182842718151806255970339330695088203366400
n[00] = 724890234288883807133692739554181389269667712903527002964305509138092747765250640043000
n[01] = 244047318040953808384618009886615469343407582143500678251608196392283700788171656780960
n[02] = 200583767506299883398799846615839618722912096662345681850356525680047866040867694686370
n[03] = 95072474317732880408873060015477036645940382387376181430971147912620913985476642401730
n[04] = -16547648260093571352064409821528871420122608189302225049240635263600825946450191092942
n[05] = -6865768954945182165273746202102702103834182475211343118142113409291904225223532915282
n[06] = 125097238907996793280228078221895618977339888699390549658898199910084249082893799887550
n[07] = 86899181507597818581114096054078208079345776935909812728525823282898418409948497684310
n[08] = 56847921916204202382738032957377262536116924684885019561903539676492216748841247695345
n[09] = 24547276233504082055846686191137024757811414077289563147946293406387265474531528641755
n[10] = -4337833911022686984336314214527962189832998337927252561590333744868926352043315516740
n[11] = -2128613158185747975129993756707705020446598436320439271687906291812329854772751114454
\end{Verbatim}
\end{minipage}

\par\medskip\noindent\begin{minipage}{\textwidth}
\textbf{$h_{36}$}
\begin{Verbatim}[fontsize=\fontsize{8}{9.4}\selectfont]
d = 2733721037906377654297837130744930514829792143182842718151806255970339330695088203366400
n[00] = 243724321419008311235385953077050887168464483820632996153070523023829526776764263685100
n[01] = 199146037199101695294576522333348032267329654498498501534842764254632528988952757336300
n[02] = 99552237782962777945734446575668234192316147426865165066142025736453823983509995947400
n[03] = 41834960618374495916923448606135598267214692703656650871110670698147980108871831462700
n[04] = -7521987311761556960596494538087691612338364731092313670582138591537998935526070960760
n[05] = -3738239736074585498384047768421480280712283333802187493597341120949346362748840653136
n[06] = 171819277401554044884462553306513336223771118192081862103303679190458550351539572469550
n[07] = 28382680992492032227575868910420724874864895485832845639664530387481429244630727082770
n[08] = 22718160367973672521699543570581690080529234476881013406836945389499578569514041890505
n[09] = 11408442049685906353588991812702818440997204486528343055042288115232218634879338057105
n[10] = -2472727747014319016917806964329164875742922562499039620066170921204923787024848410216
n[11] = -783120868471591524536529034452981675286819806984698515943462603720103226019315517636
\end{Verbatim}
\end{minipage}

\par\medskip\noindent\begin{minipage}{\textwidth}
\textbf{$h_{37}$}
\begin{Verbatim}[fontsize=\fontsize{8}{9.4}\selectfont]
d = 2733721037906377654297837130744930514829792143182842718151806255970339330695088203366400
n[00] = -659188627283393470512733592499641906893164294771735430705696036614546822723144701270300
n[01] = -259816354186404115859253078296752440755848093868815214321588302892663690875074538393260
n[02] = -191920939487370320371489209019022172362102112967996585430210831768122114773561538100820
n[03] = -91802805894742434779527963429813807393066687458554366569460166694873153403719189608840
n[04] = 15853136588607838984393618921762867076111076420295619602206316508024707173268052486152
n[05] = 6623424020982660277730332140792015589887488810705752924452205083895036950582903034424
n[06] = -145950948389173562104132836109465659836888028605918061856359150113004781191759697996150
n[07] = -81151669743461939948455863961458246302868479179168668437820775234253671339549240852210
n[08] = -59630391970241955172230381490731748445658877169359429225009077650236727470371426246465
n[09] = -25180777496992795375231065818777820357882199977503636714867896861643907090539943585405
n[10] = 4549894098731050180205023458057994500967369588491081787207102240086730571663703467830
n[11] = 2179125837240238013749139260713042491567159230439507073333272015833025166955934105010
\end{Verbatim}
\end{minipage}

\par\medskip\noindent\begin{minipage}{\textwidth}
\textbf{$h_{38}$}
\begin{Verbatim}[fontsize=\fontsize{8}{9.4}\selectfont]
d = 1366860518953188827148918565372465257414896071591421359075903127985169665347544101683200
n[00] = -405310113917233023155099066759884494882095383892066380374158646922700505342214029362300
n[01] = -224688796881114482374122011807054571252704837942491714491694377674215163995276660064020
n[02] = -163473979699209184058490595956602050811094086877074239024149922896029765543020219036920
n[03] = -74764761002985074286621610159799346816226295310098389653631939482447867722199748580760
n[04] = 12128329895439324699579033078429984149039491466884409637883426195825322053199761613394
n[05] = 5763381280089287889010554491825525247290518620235318204143006907127508535690335320842
n[06] = -132768192577360077788703600768295986838194075814118965182734388032934676801232248490550
n[07] = -51158094977006353472143650987227517450020073000833159365822500429898108572652125710990
n[08] = -41368961135045355457012303858045655715565491715515084102446608060521343904598102438915
n[09] = -18960328801446196058826174477909693897513443487320703809362756057796218313686981726905
n[10] = 3573080348584702849245324517958032122791617496990149580981899357867351014753149494456
n[11] = 1482827590440617989559770856700765005747159235907346619455537039959534126503513774916
\end{Verbatim}
\end{minipage}

\par\medskip\noindent\begin{minipage}{\textwidth}
\textbf{$h_{39}$}
\begin{Verbatim}[fontsize=\fontsize{8}{9.4}\selectfont]
d = 683430259476594413574459282686232628707448035795710679537951563992584832673772050841600
n[00] = 392663120801943552400896901991220719977390108731924233052889970467736290555117781958000
n[01] = 202045504877653826245704078244601253809701559157461368006861870684162863481805673870960
n[02] = 173185154483668746276181970756186079947428477596572793950067573631171625861338026066750
n[03] = 76948919657316977313367689998831652795894243016658504639882861733206992117144340039370
n[04] = -13605535734670015213895172119244421558827243279858456462998045889027072666316395979092
n[05] = -6067854269979903276078973862181132061117673023085767664554193459298329711890310045796
n[06] = 152003786003313447808793420385613513315762402568179271373568283548903999748792021314300
n[07] = 41080346192109271184710294659061440275348863412584052603094052612152746487579243195300
n[08] = 53742953141309790077928198076107108137353139255578571396508577600613548805649853610700
n[09] = 21857343706088014070490496918952873721913879266235851169180725805617860633164647504660
n[10] = -3919508062073923341517111566683315382619502949671728208536392467092888700431970178443
n[11] = -1703758227107723189941572893055163742852063087979235472354013635557531142341947506755
\end{Verbatim}
\end{minipage}

\par\medskip\noindent\begin{minipage}{\textwidth}
\textbf{$h_{40}$}
\begin{Verbatim}[fontsize=\fontsize{8}{9.4}\selectfont]
d = 2733721037906377654297837130744930514829792143182842718151806255970339330695088203366400
n[00] = 784060516836836547859292687798793075468508253100465294911311626203005555432155004159200
n[01] = 355545506877693945495252699684782919421066727691160366523980181823526773144013902221280
n[02] = 357351689009563224248683561165093766780761173691304773517411465560615128255193759978980
n[03] = 150712033782116927109033582777496404684627330145409705400195288941653407985661335039800
n[04] = -26582595999417388810255691599362378439903355443266375965783351962272873318858105984126
n[05] = -12428210033891608060524916593413705862420745920719537998064680227347369798011158765474
n[06] = 241402359082445377828009554152812752956942936649267976308923391849908047743738190302050
n[07] = 65478595920401337424553596239559713614334162588179660803874062777685828564350004364370
n[08] = 79027938473861341699288276193031769084042060386287397049950498289625918876291321629740
n[09] = 35485842819805246584472932764701355764720867082087292555815880038861723992278052938660
n[10] = -6192617036436452593824396039518062498389658866183469284605960400223001441979381981387
n[11] = -2526056175273302533412734508867776741195089466024855315746317446600406443753819943077
\end{Verbatim}
\end{minipage}

\par\medskip\noindent\begin{minipage}{\textwidth}
\textbf{$h_{41}$}
\begin{Verbatim}[fontsize=\fontsize{8}{9.4}\selectfont]
d = 2733721037906377654297837130744930514829792143182842718151806255970339330695088203366400
n[00] = -867460985172622450457702135135638130324762207007270910125532536011042184895666522553800
n[01] = -292620287327156394442613336186684232894186584570672389250923818638318772711787021338160
n[02] = -303929798663737575531101176676386565636200261939825104139285049216250647638106021651220
n[03] = -133274626142982485199378168408754676601829588309386137990449879558470977128204362265620
n[04] = 24102135983087619255725007853507332438868884223381859115728105358845357634688026522862
n[05] = 10576458973711931928171298503172849961093322543667205982816712464177049901723195760198
n[06] = -190431146214427481150390955543135682430150888795121278226181830322708028261117512487100
n[07] = -107276386321317995966276983241429200682171285896767563735490491401800413260442434398740
n[08] = -84491179872614936301522905856482283926798791195047146721689193688286702902462206848175
n[09] = -35308415884570753408879761185014498391484671739208330847129926193804429562312945768975
n[10] = 6003043915001389686424358246179572909778192849235172961508178166250153611618277055418
n[11] = 2874205971376045673249873058714534637886832279103525517937531893559203124662591493880
\end{Verbatim}
\end{minipage}


\clearpage
\section{Rational root boxes and positivity enclosures}\label{app:intervals}
Put $N=2^{256}$. The embedding boxes used for the exact interval checks are
\begin{equation}
 \mu\in\left[\frac{M}{10N},\frac{M+1}{10N}\right],\quad
 s\in\left[\frac{S}{N},\frac{S+1}{N}\right],\quad
 u\in\left[\frac{U}{N},\frac{U+2}{N}\right],
\end{equation}
where the integer constants are
\begin{Verbatim}[fontsize=\footnotesize]
M = 3491954164956306431453644609980892434381744577868569134761382324120276943749668
S = 258918982791360791136623237694189224638272597174970463806124557733925595847978
U = 440499284018317947084543798940668032418511528266755852350232571514445862522982
\end{Verbatim}
The sign tests described in Section~\ref{sec:positive} verify these boxes
directly. The enclosures in Table~\ref{tab:minors} are deliberately rounded
outward to simple rational endpoints: each listed integer pair $(l,r)$ means
$l/10^{15}<\Delta_{kj}<r/10^{15}$. The strict inequalities follow from the
stronger exact rational interval endpoints before outward rounding.
\begin{longtable}{@{}rrrr@{}}
\caption{Positive leading-minor enclosures, with common denominator $10^{15}$.}\label{tab:minors}\\
\toprule $k$&$j$&$l$&$r$\\\midrule\endfirsthead
\toprule $k$&$j$&$l$&$r$\\\midrule\endhead
1 & 1 & 1913474967234 & 1913474967235 \\
1 & 2 & 3147858252 & 3147858253 \\
6 & 1 & 116844895528629 & 116844895528630 \\
6 & 2 & 692437763721 & 692437763722 \\
6 & 3 & 1179239858 & 1179239859 \\
6 & 4 & 1491101 & 1491102 \\
7 & 1 & 1391263237621 & 1391263237622 \\
7 & 2 & 4228271467 & 4228271468 \\
7 & 3 & 6458386 & 6458387 \\
8 & 1 & 4585428861901 & 4585428861902 \\
8 & 2 & 14054444762 & 14054444763 \\
8 & 3 & 6458386 & 6458387 \\
9 & 1 & 5319184481121 & 5319184481122 \\
9 & 2 & 1906867713 & 1906867714 \\
\bottomrule
\end{longtable}

All these determinants are positive, establishing positive definiteness
without first assuming the validity of an LDL division. Their ratios give
the 14 positive pivots. The large integer records in the preceding appendix
can thus be interpreted as a positive operator certificate, rather than
merely a formal identity over a quotient ring.
\end{document}
